\chapter*{Preface}

OctaC is a statically typed, C-structured language for numerical and
matrix computation. Dynamically typed languages like Octave
\cite{octave-manual} let a programmer multiply two matrices of incompatible
shapes without complaint until the operation actually runs. A statically
typed language can catch that mistake before a single instruction executes,
but only if the compiler is taught to treat shape, not just type, as a
property of a value that can be checked at compile time. OctaC brings that
discipline to a C-structured surface syntax, with scalar types (\code{i16},
\code{i32}, \code{i64}, \code{f16}, \code{f32}, \code{f64}, \code{bool})
alongside matrix and vector types parameterized by their element type
(\code{matrix<i32>}, \code{vector<f32>}).

This is a living book. It begins with an introduction to OctaC and its
lexical analysis, and is meant to grow chapter by chapter as the compiler is
built over time: the formal grammar, the parser, the semantic analysis pass,
and code generation will each be added as their own chapter once they exist.
Chapter~1 introduces the language as it currently stands, and Chapter~2
covers its lexer.

This book prioritizes clarity and feasibility over completeness. Where we
simplify (omit a feature, restrict a rule, choose one semantic
over another), we try to say so explicitly rather than let the omission
speak for itself.

We hope it serves as both a reference for the finished language and an
honest record of the decisions, and trade-offs, behind it.

\begin{lstlisting}[breaklines=true]
function main() -> i32 {
    printf("Hello, world!\n");
    return 0;
}
\end{lstlisting}
